\documentclass[a4paper,10pt]{article}
\usepackage[utf8]{inputenc}
\usepackage[T1]{fontenc}
\usepackage{microtype}
\usepackage[german]{babel}
\usepackage{geometry}
\usepackage{amssymb}
\usepackage{array}
\usepackage{ragged2e}
\usepackage{graphicx}
\usepackage{parskip}
\usepackage{titlesec}

\geometry{a4paper, margin=1.6cm, top=1.8cm}
\setlength{\tabcolsep}{4pt}
\setlength{\arrayrulewidth}{0.5pt}
\renewcommand{\arraystretch}{1.75}

\newcolumntype{Q}{>{\RaggedRight\arraybackslash}m{10.2cm}}
\newcolumntype{A}{>{\centering\arraybackslash}m{1.15cm}}

\newcommand{\cb}{{\Large$\square$}}
\newcommand{\rh}[1]{\rotatebox{90}{\small\textbf{#1}\hspace{0.4em}}}

\titleformat{\section}[block]{\large\bfseries}{}{0em}{}[\vspace{-0.5em}]
\titlespacing*{\section}{0pt}{1.2em}{0.8em}

\newcommand{\tableheader}{%
  \textbf{Aussage} &
  \rh{Stimme gar nicht zu} &
  \rh{Stimme nicht zu} &
  \rh{Neutral} &
  \rh{Stimme zu} &
  \rh{Stimme sehr zu} \\
  \hline
}

\begin{document}

\begin{center}
  {\LARGE\bfseries Evaluationsbogen}\\[0.5em]
  {\small Name: \underline{\hspace{4cm}}\quad
          Datum: \underline{\hspace{3cm}}}
\end{center}

\vspace{0.5em}
{\small\textit{Bitte kreuzen Sie bei jeder Aussage die zutreffende Antwort an.}}

\section*{1. User Experience}

\noindent
\begin{tabular}{|Q|A|A|A|A|A|}
\hline
\tableheader
Ich würde das Spiel erneut spielen, um meine Bestzeit zu verbessern
  & \cb & \cb & \cb & \cb & \cb \\ \hline
Die Steuerung war intuitiv und leicht zu verstehen
  & \cb & \cb & \cb & \cb & \cb \\ \hline
Das Klettern war direkt verständlich ohne lange Erklärung
  & \cb & \cb & \cb & \cb & \cb \\ \hline
Das Spiel war abwechslungsreich (verschiedene Wege / Techniken möglich)
  & \cb & \cb & \cb & \cb & \cb \\ \hline
Ich würde das Spiel weiterempfehlen
  & \cb & \cb & \cb & \cb & \cb \\ \hline
Das Spiel hat mir Spaß gemacht
  & \cb & \cb & \cb & \cb & \cb \\ \hline
Ich hatte das Gefühl, das Movement gut zu beherrschen
  & \cb & \cb & \cb & \cb & \cb \\ \hline
\end{tabular}

\vspace{1.2em}

\section*{2. Immersion}

\noindent
\begin{tabular}{|Q|A|A|A|A|A|}
\hline
\tableheader
Ich hatte das Gefühl, wirklich im virtuellen Raum zu sein
  & \cb & \cb & \cb & \cb & \cb \\ \hline
Ich habe mich vollständig im Spiel verloren (Flow / Presence)
  & \cb & \cb & \cb & \cb & \cb \\ \hline
Das Klettern hat ein starkes Raumgefühl vermittelt
  & \cb & \cb & \cb & \cb & \cb \\ \hline
Ich habe oft vergessen, dass ich in einer realen Umgebung stehe
  & \cb & \cb & \cb & \cb & \cb \\ \hline
Geschwindigkeit und Bewegung haben die Immersion verstärkt
  & \cb & \cb & \cb & \cb & \cb \\ \hline
Ich hatte das Gefühl, wirklich an einem Berg zu hängen
  & \cb & \cb & \cb & \cb & \cb \\ \hline
Ich hatte Angst runterzufallen
  & \cb & \cb & \cb & \cb & \cb \\ \hline
\end{tabular}

\vspace{1.2em}

\section*{3. VR-Sickness / Komfort}

\noindent
\begin{tabular}{|Q|A|A|A|A|A|}
\hline
\tableheader
Ich hatte während des Spielens Kopfschmerzen
  & \cb & \cb & \cb & \cb & \cb \\ \hline
Mir wurde während des Spielens schwindelig
  & \cb & \cb & \cb & \cb & \cb \\ \hline
Ich fühlte mich nach dem Spielen unwohl
  & \cb & \cb & \cb & \cb & \cb \\ \hline
Ich musste wegen körperlicher Beschwerden eine Pause machen
  & \cb & \cb & \cb & \cb & \cb \\ \hline
Ich hatte während des Spielens Höhenangst
  & \cb & \cb & \cb & \cb & \cb \\ \hline
Ich hatte keinerlei VR-bedingte Beschwerden
  & \cb & \cb & \cb & \cb & \cb \\ \hline
\end{tabular}

\vspace{0.5em}

\section*{Anmerkungen / Feedback}

\noindent\underline{\hspace{\textwidth}}\\[1.6em]
\noindent\underline{\hspace{\textwidth}}\\[1.6em]
\noindent\underline{\hspace{\textwidth}}

\end{document}
